\documentclass[10pt,a4paper]{amsart}
\usepackage{amsmath,amssymb,amsthm,graphicx}
% Retained preamble for the Stacks Project chapter morphisms.tex (Morphisms of Schemes).
% The chapter file pinned-original-source.tex carries only \input{preamble} on its first line, so no
% preamble text lives inside the chapter itself.  The Stacks Project uses one single shared preamble for
% every chapter; the snapshot stages it verbatim in sources/stacks-chow/preamble.tex.  The block below is
% copied mechanically, line for line, from that file's lines 14-18 (the Xy-pic setup) and lines 154-217
% (the native theorem declarations together with the mathematical macros).  Nothing was added, removed or
% redefined.  The wrapper class and the amsmath/amssymb/amsthm/graphicx base packages are supplied by the
% assembler, so the class/\externaldocument part of the shared preamble is not repeated here.
% For commutative diagrams we use Xy-pic
\usepackage[all]{xy}

% We use 2cell for 2-commutative diagrams.
\xyoption{2cell}
% Theorem environments.
%
\theoremstyle{plain}
\newtheorem{theorem}[subsection]{Theorem}
\newtheorem{proposition}[subsection]{Proposition}
\newtheorem{lemma}[subsection]{Lemma}

\theoremstyle{definition}
\newtheorem{definition}[subsection]{Definition}
\newtheorem{example}[subsection]{Example}
\newtheorem{exercise}[subsection]{Exercise}
\newtheorem{situation}[subsection]{Situation}

\theoremstyle{remark}
\newtheorem{remark}[subsection]{Remark}
\newtheorem{remarks}[subsection]{Remarks}

\numberwithin{equation}{subsection}

% Macros
%
\def\lim{\mathop{\mathrm{lim}}\nolimits}
\def\colim{\mathop{\mathrm{colim}}\nolimits}
\def\Spec{\mathop{\mathrm{Spec}}}
\def\Hom{\mathop{\mathrm{Hom}}\nolimits}
\def\Ext{\mathop{\mathrm{Ext}}\nolimits}
\def\SheafHom{\mathop{\mathcal{H}\!\mathit{om}}\nolimits}
\def\SheafExt{\mathop{\mathcal{E}\!\mathit{xt}}\nolimits}
\def\Sch{\mathit{Sch}}
\def\Mor{\mathop{\mathrm{Mor}}\nolimits}
\def\Ob{\mathop{\mathrm{Ob}}\nolimits}
\def\Sh{\mathop{\mathit{Sh}}\nolimits}
\def\NL{\mathop{N\!L}\nolimits}
\def\CH{\mathop{\mathrm{CH}}\nolimits}
\def\proetale{{pro\text{-}\acute{e}tale}}
\def\etale{{\acute{e}tale}}
\def\QCoh{\mathit{QCoh}}
\def\Ker{\mathop{\mathrm{Ker}}}
\def\Im{\mathop{\mathrm{Im}}}
\def\Coker{\mathop{\mathrm{Coker}}}
\def\Coim{\mathop{\mathrm{Coim}}}

% Boxtimes
%
\DeclareMathSymbol{\boxtimes}{\mathbin}{AMSa}{"02}

%
% Macros for moduli stacks/spaces
%
\def\QCohstack{\mathcal{QC}\!\mathit{oh}}
\def\Cohstack{\mathcal{C}\!\mathit{oh}}
\def\Spacesstack{\mathcal{S}\!\mathit{paces}}
\def\Quotfunctor{\mathrm{Quot}}
\def\Hilbfunctor{\mathrm{Hilb}}
\def\Curvesstack{\mathcal{C}\!\mathit{urves}}
\def\Polarizedstack{\mathcal{P}\!\mathit{olarized}}
\def\Complexesstack{\mathcal{C}\!\mathit{omplexes}}
% \Pic is the operator that assigns to X its picard group, usage \Pic(X)
% \Picardstack_{X/B} denotes the Picard stack of X over B
% \Picardfunctor_{X/B} denotes the Picard functor of X over B
\def\Pic{\mathop{\mathrm{Pic}}\nolimits}
\def\Picardstack{\mathcal{P}\!\mathit{ic}}
\def\Picardfunctor{\mathrm{Pic}}
\def\Deformationcategory{\mathcal{D}\!\mathit{ef}}

% --- appended 2026-10-12 (Chief): the project own definitions of the reference/slogan/history
% display environments, copied mechanically from the shared Stacks preamble lines 8-12
% (\usepackage{verbatim} plus the three \newenvironment lines, which the chapter preamble keeps
% verbatim from the project preamble). Nothing added, removed or redefined: these environments
% carry no mathematical content and the project itself typesets them as comments. Without these
% three lines a counted statement whose own text uses them fails with Environment slogan undefined.
\usepackage{verbatim}
\newenvironment{reference}{\comment}{\endcomment}
\newenvironment{slogan}{\comment}{\endcomment}
\newenvironment{history}{\comment}{\endcomment}
\usepackage{cleveref}
\crefname{originallemma}{original lemma}{original lemmas}
\Crefname{originallemma}{Original lemma}{Original lemmas}
\makeatletter
\newlabel{morphisms-lemma-birational-integral}{{0BAA}{source}{}{sourcepointer-1}{}}
\newlabel{morphisms-lemma-birational-integral@cref}{{[originallemma][0][]0BAA}{[1][0][]source}{}{}{}}
\makeatother
\crefname{originallemma}{original lemma}{original lemmas}
\Crefname{originallemma}{Original lemma}{Original lemmas}
\makeatletter
\newlabel{schemes-lemma-section-immersion}{{01KT}{source}{}{sourcepointer-2}{}}
\newlabel{schemes-lemma-section-immersion@cref}{{[originallemma][0][]01KT}{[1][0][]source}{}{}{}}
\makeatother
\crefname{originallemma}{original lemma}{original lemmas}
\Crefname{originallemma}{Original lemma}{Original lemmas}
\makeatletter
\newlabel{schemes-lemma-immersions-monomorphisms}{{01L7}{source}{}{sourcepointer-3}{}}
\newlabel{schemes-lemma-immersions-monomorphisms@cref}{{[originallemma][0][]01L7}{[1][0][]source}{}{}{}}
\makeatother
\title{Birationality of irreducible schemes is equivalent to having isomorphic nonempty opens}
\author{The Stacks Project Authors}
\date{stacks-project:morphisms.tex:local-snapshot-20261012; GNU FDL 1.2 or later}
\begin{document}
\maketitle
\setcounter{subsection}{1}
\noindent\textit{Editorial note.} All delivered mathematical text is the author's own, copied from the pinned Stacks Project chapter file morphisms.tex; no statement or proof step was rewritten, reordered or removed.\par\medskip
\noindent\textit{Editorial note.} The chapter file pinned-original-source.tex carries only \textbackslash{}input\{preamble\} on its first line, so no preamble text lives inside the chapter. The retained file sources/stacks-morphisms/retained-preamble.tex is a mechanical, line-for-line copy of the project's single shared preamble as staged in the snapshot. Nothing was added, removed or redefined, and the excerpt serves only as preamble, never as a span coordinate.\par\medskip
\noindent\textit{Editorial note.} Both counted spans start at the \textbackslash{}begin line of their environment, so the statement and the author's proof are delivered whole; surrounding prose of the section is deliberately outside the counted spans and is carried as standard prerequisites instead.\par\medskip
\noindent\textit{Editorial note.} The citation fields (source\_id, source\_version, authors, source\_url, license.name, source\_title) are identical to those of every other card this lane has produced from this chapter, so the batch selector treats them as one source.\par\medskip
\noindent\textit{Editorial note.} The relative (over a base scheme) part of the statement is dispatched by the author with 'In the \$S\$-rational maps case, the exact same argument works.' -- a self-described repetition of the argument just completed, accepted with the wording quoted under the Lead's 12:4x rule, not treated as an omitted statement item.\par\medskip
\noindent\textit{Allowed standard prerequisites (editorial).}\begin{itemize}
\item A section of a morphism whose composite is the identity is an immersion. Schemes, Lemma \textbackslash{}ref\{schemes-lemma-section-immersion\}, cited at proof lines 12919-12921. Shows the morphism between the chosen opens is an immersion.
\item Immersions are monomorphisms. Schemes, Lemma \textbackslash{}ref\{schemes-lemma-immersions-monomorphisms\}, cited at proof lines 12925-12927. Turns a monomorphism which is an immersion into an open immersion.
\item Dominant rational maps and their composition (the definition of birationality). Definition \textbackslash{}ref\{definition-birational\}, the definition this lemma works with. Supplies the inverse rational maps the proof starts from.
\end{itemize}\medskip
\ifdefined\hypertarget\hypertarget{sourcepointer-1}{}\fi
\noindent\textit{Source reference.} Original lemma 0BAA: focused-seeds/AGP/sources/stacks-snapshot-20261011/tags.txt, source lines 14203--14203. The lemma's own canonical tag, quoted so that readers can resolve the label inside the delivered text; it points at the counted statement itself and carries no proof content.\par\medskip
\ifdefined\hypertarget\hypertarget{sourcepointer-2}{}\fi
\noindent\textit{Source reference.} Original lemma 01KT: focused-seeds/AGP/sources/stacks-snapshot-20261011/tags.txt, source lines 2062--2062. Navigation pointer to a statement the counted proof refers to but does not contain; nothing load-bearing is delegated to it.\par\medskip
\ifdefined\hypertarget\hypertarget{sourcepointer-3}{}\fi
\noindent\textit{Source reference.} Original lemma 01L7: focused-seeds/AGP/sources/stacks-snapshot-20261011/tags.txt, source lines 2076--2076. Navigation pointer to a statement the counted proof refers to but does not contain; nothing load-bearing is delegated to it.\par\medskip
\begin{lemma}
\label{lemma-birational-integral}
Let $X$ and $Y$ be irreducible schemes.
\begin{enumerate}
\item The schemes $X$ and $Y$ are birational if and only if they have
isomorphic nonempty opens.
\item Assume $X$ and $Y$ are schemes over a base scheme $S$. Then
$X$ and $Y$ are $S$-birational if and only if there are nonempty
opens $U \subset X$ and $V \subset Y$ which are $S$-isomorphic.
\end{enumerate}
\end{lemma}

\begin{proof}
Assume $X$ and $Y$ are birational. Let $f : U \to Y$ and $g : V \to X$
define inverse dominant rational maps from $X$ to $Y$ and from $Y$ to $X$.
We may assume $V$ affine. We may replace $U$ by an affine open of $f^{-1}(V)$.
As $g \circ f$ is the identity as a dominant rational map, we see that
the composition $U \to V \to X$ is the identity on a dense open of $U$.
Thus after replacing $U$ by a smaller affine open we may assume that
$U \to V \to X$ is the inclusion of $U$ into $X$. It follows that
$U \to V$ is an immersion
(apply Schemes, Lemma \ref{schemes-lemma-section-immersion}
to $U \to g^{-1}(U) \to U$).
However, switching the roles of $U$ and $V$ and redoing the argument
above, we see that there exists a nonempty affine open $V' \subset V$
such that the inclusion factors as $V' \to U \to V$. Then $V' \to U$ is
necessarily an open immersion. Namely, $V' \to f^{-1}(V') \to V'$ are
monomorphisms
(Schemes, Lemma \ref{schemes-lemma-immersions-monomorphisms})
composing to the identity, hence isomorphisms.
Thus $V'$ is isomorphic to an open of both $X$ and $Y$.
In the $S$-rational maps case, the exact same argument works.
\end{proof}

\end{document}
